%================================================================
\chapter{Assessment and Progress Record}
\label{ch:assessment}
%================================================================

\begin{Form}

This chapter is a fillable PDF form intended for faculty and staff
who are onboarding as contributors to the Department of Mathematics
curriculum repository. Each section corresponds to one category of
readiness: Required, Essential, Desirable, and Optional. Checkboxes
record binary completion of verification steps; text fields record
short written responses that demonstrate conceptual understanding.
Use this chapter as a running checklist during onboarding.

\vs
\noindent
\textbf{Saving your responses requires Adobe Acrobat Reader} (free
for Windows, macOS, and Linux; available at
\url{https://get.adobe.com/reader/}). Other viewers behave as
follows:

\begin{itemize}
    \item \textbf{macOS Preview:} displays form fields and allows
          typing in them, but \emph{does not save the entered data}.
          If you close and reopen the file in Preview, your responses
          will be lost.
    \item \textbf{Google Chrome (built-in PDF viewer):} displays and
          accepts input in form fields, but \emph{does not save the
          data}. Printing to PDF from Chrome will not preserve field
          contents either.
    \item \textbf{Most Linux PDF viewers} (Evince, Okular, Zathura):
          \emph{ignore form fields entirely}. Fields will not appear
          interactive.
    \item \textbf{Adobe Acrobat Reader:} fully saves, prints, and
          exports form data. Use \textbf{File $>$ Save} (not Save As)
          after completing each section.
\end{itemize}

\vs
\noindent
\textbf{How to interpret the four categories:}

\begin{description}
    \item[Required] Core repository tasks. These must be completed
          before a contributor can make any changes to the curriculum
          document.
    \item[Essential] Local environment setup: \LaTeX{}, Python, Git,
          and Visual Studio Code. These must be working correctly
          before beginning any authoring or scripting work.
    \item[Desirable] Conceptual understanding of the tools introduced
          in the infrastructure chapters. Completing these
          demonstrates readiness to work independently on the
          repository.
    \item[Optional] Local LLM infrastructure (Ollama). Not required
          for curriculum authoring; relevant for contributors who
          wish to use local language models in their workflow.
\end{description}
\vs

%================================================================
\section{Required: Repository Onboarding}
\label{sec:assess-required}
%================================================================

The following items verify that you can retrieve, compile, and
contribute to the curriculum repository. All items must be
completed before editing any source file.

\vs

\noindent\CheckBox[name=r1,width=0.4cm,height=0.4cm]{}~\textbf{R1.}
\texttt{git clone} of the repository completes without errors and
the local working copy contains the \texttt{curriculum/} folder.
(See Chapter~\ref{ch:git}, Section~\ref{subsec:multiple-remotes}.)

\vs

\noindent\CheckBox[name=r2,width=0.4cm,height=0.4cm]{}~\textbf{R2.}
Running \texttt{pdflatex MathCurriculum.tex} from inside the
\texttt{curriculum/} folder produces \texttt{MathCurriculum.pdf}
with no fatal errors.
(See Chapter~\ref{ch:latex}, Section~\ref{sec:first-latex-doc}.)

\vs

\noindent\CheckBox[name=r3,width=0.4cm,height=0.4cm]{}~\textbf{R3.}
Running \texttt{gh} (or \texttt{gh.bat} on Windows) with a commit
message successfully stages, commits, and pushes a change to the
repository, and \texttt{VERSION.txt} is updated with the new hash.
(See Chapter~\ref{ch:git}.)

\vs

\noindent\CheckBox[name=r4,width=0.4cm,height=0.4cm]{}~\textbf{R4.}
The commit hash in \texttt{VERSION.txt} matches the value returned
by \texttt{git rev-parse --short HEAD} in your local repository.
(See Chapter~\ref{ch:git}.)

\vs

\noindent\textbf{R5.} Describe the two-part structure of the
curriculum document. What is contained in the infrastructure
chapters, and what is contained in the course-unit chapters? How
are unit chapters identified in the source tree?
(See the Introduction.)\\[4pt]
\TextField[name=r5,width=\linewidth,height=3cm,multiline=true,
           bordercolor=black]{}

%================================================================
\section{Essential: Environment and Toolchain Setup}
\label{sec:assess-essential}
%================================================================

Check each box when the corresponding verification step succeeds.
All items in this section should be checked before beginning any
authoring or scripting work on the repository.

\vs

\noindent\CheckBox[name=e1,width=0.4cm,height=0.4cm]{}~\textbf{E1.}
\texttt{pdflatex --version} returns a version number in the
terminal.
(See Chapter~\ref{ch:environment-setup}.)

\vs

\noindent\CheckBox[name=e2,width=0.4cm,height=0.4cm]{}~\textbf{E2.}
\texttt{python --version} returns a version number in the terminal.
(See Chapter~\ref{ch:environment-setup}, Section~\ref{sec:python}.)

\vs

\noindent\CheckBox[name=e3,width=0.4cm,height=0.4cm]{}~\textbf{E3.}
\texttt{git --version} returns a version number.
(See Chapter~\ref{ch:environment-setup}, Section~\ref{sec:git}.)

\vs

\noindent\CheckBox[name=e4,width=0.4cm,height=0.4cm]{}~\textbf{E4.}
\texttt{code .} opens Visual Studio Code from a terminal window.
(See Chapter~\ref{ch:environment-setup}, Section~\ref{sec:vscstart}.)

\vs

\noindent\CheckBox[name=e5,width=0.4cm,height=0.4cm]{}~\textbf{E5.}
The Python virtual environment activates without errors, and the
packages listed in \texttt{pyproject.toml} (or
\texttt{requirements.txt}) import cleanly in that environment.
(See Chapter~\ref{ch:python-environments},
Section~\ref{sec:project-libraries}.)

\vs

\noindent\CheckBox[name=e6,width=0.4cm,height=0.4cm]{}~\textbf{E6.}
Running \texttt{latexmk -pdf MathCurriculum.tex} produces a
correctly paginated PDF with a table of contents and all
cross-references resolved.
(See Chapter~\ref{ch:latex}.)

%================================================================
\section{Desirable: Technical Understanding}
\label{sec:assess-desirable}
%================================================================

These items demonstrate working understanding of the tools
introduced in the infrastructure chapters. Written responses need
not be long; two or three sentences each are sufficient.

\vs

\noindent\textbf{D1.} What is the current working directory when you
first open a terminal? Write the command you would use to navigate
to the \texttt{curriculum/} folder inside the cloned repository.
(See Chapter~\ref{ch:cli-scripting},
Section~\ref{sec:what-is-a-file}.)\\[4pt]
\TextField[name=d1,width=\linewidth,height=2.5cm,multiline=true,
           bordercolor=black]{}

\vs

\noindent\textbf{D2.} What is the difference between a Python
virtual environment and the global Python installation? Give one
scenario in the context of this repository where you would need
a separate virtual environment.
(See Chapter~\ref{ch:python-environments}.)\\[4pt]
\TextField[name=d2,width=\linewidth,height=2.5cm,multiline=true,
           bordercolor=black]{}

\vs

\noindent\textbf{D3.} Why does \texttt{git push} require a prior
\texttt{git commit}? What is the conceptual difference between
the two operations, and why does the \texttt{gh.bat} script
perform both in sequence?
(See Chapter~\ref{ch:git}.)\\[4pt]
\TextField[name=d3,width=\linewidth,height=2.5cm,multiline=true,
           bordercolor=black]{}

\vs

\noindent\textbf{D4.} In \LaTeX{}, what is the practical difference
between \verb|\[...\]| and \verb|$$...$$|? Which form should be
used in new curriculum source files, and why?
(See Chapter~\ref{ch:latex}, Section~\ref{sec:latex-math}.)\\[4pt]
\TextField[name=d4,width=\linewidth,height=2.5cm,multiline=true,
           bordercolor=black]{}

\vs

\noindent\textbf{D5.} What does the \texttt{regen.py} script do
when run for a given course unit? Describe the input it reads,
the output it produces, and the role of the \texttt{ASSESSMENTS}
dictionary in controlling the output.
(See the repository \texttt{Script/} folder and associated
documentation.)\\[4pt]
\TextField[name=d5,width=\linewidth,height=2.5cm,multiline=true,
           bordercolor=black]{}

\vs

\noindent\textbf{D6.} Explain the label scheme used for
cross-references in the curriculum source files. What is the
format of a section label, how is the four-character mnemonic
derived, and how is the hash suffix computed?
(See the curriculum pipeline documentation.)\\[4pt]
\TextField[name=d6,width=\linewidth,height=2.5cm,multiline=true,
           bordercolor=black]{}

%================================================================
\section{Optional: Local LLM Infrastructure (Ollama)}
\label{sec:assess-optional}
%================================================================

These items apply only to contributors who install and configure
the local Ollama LLM infrastructure. They are not required for
curriculum authoring or repository maintenance.

\vs

\noindent\CheckBox[name=o1,width=0.4cm,height=0.4cm]{}~\textbf{O1.}
Ollama is installed and \texttt{ollama list} returns at least one
model in the terminal.
(See Chapter~\ref{ch:ollama-infrastructure},
Section~\ref{sec:install}.)

\vs

\noindent\CheckBox[name=o2,width=0.4cm,height=0.4cm]{}~\textbf{O2.}
\texttt{connectivity\_test.py} passes all verification checks
(TCP handshake, model list, inference).
(See Chapter~\ref{ch:ollama-infrastructure}.)

\vs

\noindent\textbf{O3.} Explain the tradeoff between context window
size and memory consumption in Ollama. Why is
\texttt{OLLAMA\_NUM\_CTX=8192} a safe default for large models on
hardware with limited VRAM?
(See Chapter~\ref{ch:ollama-infrastructure},
Section~\ref{sec:memory}.)\\[4pt]
\TextField[name=o3,width=\linewidth,height=3cm,multiline=true,
           bordercolor=black]{}

\end{Form}